%%% ============================================================
%%% Freeze-Thaw Transit: Dual-Regime Scalar Dark Energy
%%% with Analytic Vacuum at e^{-1} from a Logarithmic Potential
%%%
%%% B.R. Sanders, M. Gish, H.J. Johnson (2026)
%%% Target venue: JCAP (Journal of Cosmology and Astroparticle Physics)
%%% MSC: 83F05, 83C56, 85A40
%%% PACS: 95.36.+x (dark energy); 98.80.Es (observational cosmology);
%%%       98.80.Cq (cosmological perturbations)
%%%
%%% Companions in the same submission cycle:
%%%   - Sanders & Gish (2026), "Non-Associativity Decay in Binary
%%%     Composition Tables over Z/NZ" -- JCT-A
%%%   - Sanders & Gish (2026), "Joint Closure, Per-Coordinate Fuse
%%%     Data, and a Closed-Form Algebraic Attractor of Two
%%%     Commutative Binary Operations on Z/10Z" -- Algebraic
%%%     Combinatorics
%%%
%%% DOI for verification scripts: 10.5281/zenodo.18852047
%%% ============================================================

\documentclass[11pt,reqno]{amsart}

\usepackage{amsmath,amssymb,amsthm,mathtools}
\usepackage[margin=1in]{geometry}
\usepackage[colorlinks=true,linkcolor=blue,citecolor=blue,urlcolor=blue]{hyperref}
\usepackage{microtype}
\usepackage{array}

%% -------- theorem environments --------
\theoremstyle{plain}
\newtheorem{theorem}{Theorem}
\newtheorem{proposition}[theorem]{Proposition}
\newtheorem{corollary}[theorem]{Corollary}
\newtheorem{lemma}[theorem]{Lemma}

\theoremstyle{definition}
\newtheorem{definition}[theorem]{Definition}

\theoremstyle{remark}
\newtheorem*{remark}{Remark}

%% -------- macros --------
\newcommand{\Pl}{\mathrm{Pl}}
\newcommand{\SM}{\mathrm{SM}}
\newcommand{\Gibbs}{\mathrm{Gibbs}}
\DeclareMathOperator{\diag}{diag}

%% -------- title matter --------
\title[Freeze-Thaw Transit]{Freeze-Thaw Transit: Dual-Regime\\
Scalar Dark Energy with Analytic Vacuum at $e^{-1}$\\
from a Logarithmic Potential}

\author{B.~R.~Sanders}
\address{7Site LLC, Hot Springs, Arkansas, USA}
\email{brayden@7site.co}

\author{M.~Gish}
\address{Independent Researcher, Hot Springs, Arkansas, USA}
\email{monica.gish1992@gmail.com}

\author{H.~J.~Johnson}
\address{Independent Researcher, Billings, Montana, USA}
\email{hjj01986@gmail.com}

\date{\today}

\keywords{quintessence, dark energy, scalar field, logarithmic potential,
freeze-thaw transit, dual-regime quintessence, equation of state, DESI,
information-theoretic cosmology}

\subjclass[2020]{83F05, 83C56, 85A40}

\begin{document}

\begin{abstract}
We study a quintessence dark-energy model in which a real, positive,
dimensionless scalar field $\Xi(x)$ is minimally coupled to gravity with
self-interaction $V(\Xi) = \Lambda^{4}\,\Xi\log\Xi$, where $\Lambda$ is a
mass scale that sets the potential energy density. The potential admits an
analytic minimum at $\Xi_0 = e^{-1}$. As is generic for any dimensionless
field whose coupling enters as an overall prefactor, the location of this
minimum is independent of $\Lambda$; what is specific to the choice
$V \propto \Xi\log\Xi$ is the value $\Xi_0 = e^{-1}$ itself. Expanding about
$\Xi_0$ gives a stable massive scalar with
$m_\Xi^{2} = \Lambda^{4} e / M_\Pl^{2}$, which places $\Lambda$ numerically
near the observed dark-energy scale when $m_\Xi \sim H_0$, giving
$\Lambda \approx 1.7$~meV. In a spatially flat FRW background with
outbound initial condition $\dot\Xi_i > 0$ at $z_i \approx 20$, the field
equation drives a \emph{dual-regime} cosmological history we term
\emph{freeze-thaw transit}: the field thaws outbound (Type-T), reaches
an instantaneous frozen turnaround at intermediate redshift $z_\star
\approx 2$ where $\dot\Xi = 0$ and $w_\Xi(z_\star) = -1$ momentarily
(Type-F), and then asymptotically refreezes back toward the vacuum
$\Xi_0 = e^{-1}$ as $z \to -1$ (Type-A), with $w \to -1$ at the late-time
endpoint. This trajectory traverses both standard quintessence regimes
(thawing and freezing in the Caldwell-Linder \cite{CaldwellLinder2005}
sense) within a single physical history; its observational signature is
a non-monotone $w_{\rm DE}(z)$ with a local minimum near $-1$ at
intermediate redshift, formalized as falsification criterion $(F_5)$
below. We provide an illustrative background-level numerical integration
showing that the model can reproduce $(w_0, w_a)$ values close to the
DESI~2024 DR1 published Gaussian summary on the CPL parametrization for
appropriately chosen initial conditions; the $\chi^{2}$ values reported
here quantify proximity to the published $(w_0, w_a)$ marginal Gaussian
and are not derived from the underlying joint BAO + CMB + SN likelihood,
which is deferred to a companion paper. The minimal action couples to
matter only gravitationally, so no observable fifth force is expected
from the minimal theory. The functional form $-\Xi\log\Xi$ coincides
with the per-bin Gibbs integrand on a normalized probability
distribution and with the Bialynicki-Birula--Mycielski nonlinearity in
nonlinear quantum mechanics \cite{BialynickiBirulaMycielski1976}; both
connections are structural rhymes (the present $\Xi$ is unbounded and
unnormalized, not a probability density), cited as motivation for
studying this functional form rather than as derivations.
Falsifiability is anchored in five Stage-IV predictions: rolling-branch
$w_\Xi(z) \geq -1$, monotone or non-monotone shape per dual-regime
hypothesis, asymptotic limit $w \to -1$, two-parameter $w(z)$ profile,
and the local-minimum signature of the Type-F turnaround.
\end{abstract}

\maketitle

\section{Introduction}

The nature of dark energy is the central open question in observational
cosmology \cite{SahniStarobinsky2000,Padmanabhan2003,PeeblesRatra2003,FriemanTurnerHuterer2008,JoyceJainKhouryTrodden2015}.
Recent results from the Dark Energy Spectroscopic Instrument
\cite{DESI2024BAO,DESI2024VI} suggest a preference for dynamical dark energy
over the cosmological constant at the $2.8$--$4.2\sigma$ level in joint
analyses combining DESI BAO with CMB and supernova data, in the
Chevallier-Polarski-Linder \cite{CPL2001,Linder2003} parametrization
$w(z) = w_0 + w_a z/(1+z)$. The DR1 BAO data alone are consistent with flat
$\Lambda$CDM; the dynamical-dark-energy preference appears in combined
analyses. These results have revived interest in scalar-field models of
quintessence \cite{RatraPeebles1988,Wetterich1988,Frieman1995,CaldwellDaveSteinhardt1998,SteinhardtWangZlatev1999,ZlatevWangSteinhardt1999}.
Standard quintessence potentials --- inverse power laws, exponentials,
pseudo-Nambu-Goldstone bosons --- are motivated by high-energy physics
(supersymmetry breaking, axion physics, string moduli) and share the
generic property that the field rolls toward infinity or has no
well-defined late-time attractor distinct from zero. The dynamical
properties of these models, including tracker behavior, freezing and
thawing classification, and dependence on initial conditions, have been
extensively reviewed \cite{CopelandSamiTsujikawa2006,Tsujikawa2013}.

We study a specific quintessence potential:
\begin{equation}\label{eq:Vxi}
  V(\Xi) = \Lambda^{4}\,\Xi\log\Xi,
  \qquad \Xi > 0,\ \Lambda > 0,
\end{equation}
where $\Xi$ is a real, positive, dimensionless scalar field and $\Lambda$ is
a mass scale that sets the potential energy density. We shall call this the
\emph{logarithmic quintessence potential}. It has an analytic minimum at
$\Xi_0 = e^{-1}$. Because $\Xi$ is dimensionless and $\Lambda^{4}$ enters
as an overall prefactor, the location of this minimum is independent of
$\Lambda$; this is generic for any dimensionless-field potential of the
form $\Lambda^{4} f(\Xi)$ and is not a feature peculiar to $\Xi\log\Xi$.
The feature peculiar to this choice is the value $\Xi_0 = e^{-1}$, which
is also the maximizer of the per-bin Gibbs integrand $-x\log x$ over
$x \in (0,1]$. We discuss this structural connection in
\S\ref{sec:entropy} but do not treat it as derivational: $\Xi$ is an
unbounded real scalar field rather than a normalized probability density,
so the connection to thermodynamic entropy is a structural rhyme rather
than an operational identification.

The substantive cosmological observation of the present paper is that on
the physical rolling branch of \eqref{eq:Vxi} with outbound initial
condition $\dot\Xi_i > 0$, the FRW trajectory is \emph{dual-regime}: the
field thaws outbound from $\Xi_i$ near the vacuum, reaches an
instantaneous frozen turnaround at intermediate redshift $z_\star
\approx 2$ where $\dot\Xi = 0$ momentarily and the bare scalar EoS
$w_\Xi(z_\star) = -1$, then asymptotically refreezes back toward
$\Xi_0 = e^{-1}$ as $z \to -1$. We classify these three regimes as
Type-T (thawing), Type-F (frozen turnaround), and Type-A (asymptotic
refreeze) in \S\ref{sec:typeTFA}. Caldwell-Linder
\cite{CaldwellLinder2005} classify standard quintessence into two
mutually exclusive classes: \emph{freezing}, with monotone $w(z) \to -1$
from above, and \emph{thawing}, with monotone departure from $-1$ over
the observable epoch. The dual-regime trajectory studied here is in
neither class --- it traverses both within a single physical history
--- and its observational signature is a non-monotone $w_{\rm DE}(z)$
with a local minimum near $-1$ at intermediate redshift. This signature
is not produced by any single-regime quintessence and is formalized as
falsification criterion $(F_5)$ in \S\ref{sec:falsify}. The
present-epoch value $w_\Xi(z=0) \approx -0.79$ in the documented fit
\eqref{eq:desibestfit} reflects the field's position on the inbound leg
of the trajectory, having passed through the Type-F turnaround near
$z \approx 2$.

\paragraph{Companions in the synthesis program.}
Two companion papers in the same submission cycle place this
cosmological model in a broader algebraic-combinatorial setting.
Sanders \& Gish (2026), \emph{Non-Associativity Decay in Binary
Composition Tables over $\mathbb{Z}/N\mathbb{Z}$}
\cite{SandersGishSigma2026}, target \emph{Journal of Combinatorial
Theory, Series A}, proves a $\sigma(N) \leq 2/N$ rate theorem for
non-associativity decay on a specific family of squarefree-modular
composition tables; the rate vanishes in the continuum limit and is the
combinatorial-side input that, combined with the Bialynicki-Birula
separability uniqueness theorem
\cite{BialynickiBirulaMycielski1976}, motivates the logarithmic
functional form $\Xi\log\Xi$ studied here as a candidate quintessence
potential. Sanders \& Gish (2026), \emph{Joint Closure,
Per-Coordinate Fuse Data, and a Closed-Form Algebraic Attractor of Two
Commutative Binary Operations on $\mathbb{Z}/10\mathbb{Z}$}
\cite{SandersGishFourCore2026}, target \emph{Algebraic Combinatorics},
extracts an explicit closed-form algebraic structure
($h/\beta = 1+\sqrt{3}$ at the symmetric mixing weight $\alpha = 1/2$,
governed by the integer quartic with Galois group $D_4$ and number
field LMFDB~$4.2.10224.1$) from the same composition-table family. We
emphasize at the outset that the present paper does \emph{not} derive
the action \eqref{eq:action} from the companions; the structural
relations between the substrate combinatorics and the cosmological
potential are motivational, and the empirical content of
\S\ref{sec:obs} stands or falls on its own. Forward synthesis to a
unified framework is the subject of separate work in preparation.

The paper is organized as follows. In \S\ref{sec:action} we state the
canonical action, fix unit conventions, and delimit the physical phase space.
In \S\ref{sec:derivations} we derive the Euler-Lagrange equation, the
stress-energy tensor, the analytic vacuum $\Xi_0 = e^{-1}$, the fluctuation
mass $m_\Xi^{2} = \Lambda^{4} e / M_\Pl^{2}$, and the bounded-below
global behavior of $V$.
Section \S\ref{sec:frw} specializes to a spatially flat FRW background and
derives the field energy density, pressure, and equation-of-state function.
Section \S\ref{sec:entropy} discusses the structural connection to the
Gibbs functional and to the Bialynicki-Birula--Mycielski nonlinearity in
nonlinear quantum mechanics. Section \S\ref{sec:obs} presents the
observational program, including an illustrative consistency check against
the DESI~2024 DR1 $(w_0, w_a)$ Gaussian summary, an explicit acknowledgment
of the methodological scope of that check, and discussions of EFT validity
over the field excursion and of the model's perturbation properties.
Section \S\ref{sec:comparison} compares the model to standard quintessence
potentials and to logarithmic prior art in the dark-sector literature.
Section \S\ref{sec:scope} states what the theory does \emph{not} claim.

\section{Canonical Action}\label{sec:action}

\subsection{Field content and conventions}\label{sec:fieldcontent}

The dynamical field content is
\[
  \bigl\{ g_{\mu\nu},\ \text{SM fields},\ \Xi(x) \bigr\},
\]
where $g_{\mu\nu}$ is a Lorentzian metric in the $(-,+,+,+)$ signature
convention, the SM sector follows the usual Standard Model prescription,
and $\Xi(x)$ is a real, positive, dimensionless scalar. $\Xi$ is a gauge
singlet and carries no charge under any internal symmetry, so the
minimal action contains no associated conserved current.

The effective theory is defined on the field domain $\Xi > 0$. In the
numerical analysis of \S\ref{sec:desifit} we restrict attention to initial
data $(\Xi_i, \dot\Xi_i)$ at $z_i \approx 20$ whose FRW evolution remains
within this domain over the redshift interval $0 \leq z \leq z_i$;
configurations leading to $\Xi \to 0^{+}$ during cosmic evolution are
outside the field domain considered in this effective theory and are
not considered.\footnote{An
equivalent formulation reparametrizes via $\psi = \log\Xi \in \mathbb{R}$,
eliminating the positivity constraint at the cost of a $\Xi$-dependent
kinetic prefactor $\propto e^{2\psi}$. A canonical dimension-1 field
$\phi = M_\Pl \Xi$ can also be introduced; in those variables the kinetic
term takes the standard form $\tfrac{1}{2}(\partial\phi)^{2}$ and the
potential becomes $V(\phi) = \Lambda^{4}(\phi/M_\Pl) \log(\phi/M_\Pl)$.
We work with the dimensionless field $\Xi$ throughout because in these
variables the stationary point $\Xi_0 = e^{-1}$ is especially transparent.}

\subsection{Action}

The minimal action is
\begin{equation}\label{eq:action}
  S = \int d^{4}x\,\sqrt{-g}\left[\frac{R}{16\pi G}
     + \mathcal{L}_{\SM}
     - \frac{M_\Pl^{2}}{2}\,g^{\mu\nu}\partial_\mu\Xi\,\partial_\nu\Xi
     - \Lambda^{4}\,\Xi\log\Xi\right],
\end{equation}
where $M_\Pl^{2} = (8\pi G)^{-1}$ is the reduced Planck mass and
$\Lambda$ is a mass scale that sets the potential energy density. The
signs of the kinetic and potential terms in \eqref{eq:action} are those
required by the $(-,+,+,+)$ signature convention adopted in
\S\ref{sec:fieldcontent}: for a homogeneous configuration $\Xi(t)$,
$g^{\mu\nu}\partial_\mu\Xi\,\partial_\nu\Xi = -\dot\Xi^{2}$, and the
displayed bracket evaluates to $\tfrac{1}{2}M_\Pl^{2}\dot\Xi^{2} -
V(\Xi)$, which is the standard non-ghost canonical form. Both
$\Xi$-dependent terms in the bracket have mass dimension 4: the kinetic
term acquires dimension $[M^{4}]$ from the $M_\Pl^{2}$ prefactor and the
two derivatives acting on the dimensionless $\Xi$, and the potential
term acquires dimension $[M^{4}]$ from the $\Lambda^{4}$ prefactor
multiplying the dimensionless combination $\Xi\log\Xi$. The single
model-specific parameter is $\Lambda$ (or equivalently the dimensionless
ratio $\Lambda^{4}/M_\Pl^{4}$); $\Lambda$ is fixed phenomenologically by
demanding the model reproduce the observed dark-energy density today,
which (as we show in \S\ref{sec:derivations}) places $\Lambda$
numerically near the dark-energy scale.

\begin{remark}
The model derives its predictions from $\Xi$ and $\Lambda$ alone; any
information-theoretic or entropic origin story for the potential form is
supplementary to the action and is deferred to \S\ref{sec:entropy}.
\end{remark}

\section{Derivations}\label{sec:derivations}

\subsection{Euler-Lagrange equation}

From $\mathcal{L}_\Xi = -\tfrac{1}{2} M_\Pl^{2}\, g^{\mu\nu}\partial_\mu\Xi\,\partial_\nu\Xi
- \Lambda^{4}\,\Xi\log\Xi$,
\[
  \frac{\partial \mathcal{L}_\Xi}{\partial \Xi} = -\Lambda^{4}(1 + \log\Xi),
  \qquad
  \nabla_\mu \frac{\partial \mathcal{L}_\Xi}{\partial(\nabla_\mu \Xi)}
  = -M_\Pl^{2}\,\Box\Xi.
\]
The Euler-Lagrange equation
$\partial\mathcal{L}_\Xi/\partial\Xi = \nabla_\mu(\partial\mathcal{L}_\Xi/\partial(\nabla_\mu\Xi))$
gives
\begin{equation}\label{eq:EL}
  \boxed{\;M_\Pl^{2}\,\Box\Xi = \Lambda^{4}(1 + \log\Xi)\;}
\end{equation}
or equivalently $\Box\Xi = (\Lambda^{4}/M_\Pl^{2})(1+\log\Xi)$, which
exhibits the characteristic mass scale $\Lambda^{4}/M_\Pl^{2}$ that
governs the dynamics. The vacuum condition $\Box\Xi = 0$ gives
$1 + \log\Xi_0 = 0$ as a coupling-independent algebraic equation, with
the prefactor $\Lambda^{4}/M_\Pl^{2}$ canceling at the vacuum.

\subsection{Stress-energy tensor}

Varying \eqref{eq:action} with respect to $g^{\mu\nu}$ gives
\begin{equation}\label{eq:stress}
  T^{\Xi}_{\mu\nu} = M_\Pl^{2}\,\partial_\mu\Xi\,\partial_\nu\Xi
  - g_{\mu\nu}\!\left[\frac{M_\Pl^{2}}{2}\,g^{\alpha\beta}\partial_\alpha\Xi\,\partial_\beta\Xi
  + \Lambda^{4}\,\Xi\log\Xi\right],
\end{equation}
manifestly symmetric in $(\mu, \nu)$. Covariant conservation
$\nabla^{\mu} T^{\Xi}_{\mu\nu} = 0$ holds on shell as a direct consequence
of \eqref{eq:EL} together with the contracted Bianchi identity, so
\eqref{eq:stress} is a consistent source for the Einstein equation
$G_{\mu\nu} = 8\pi G\, T^{\SM}_{\mu\nu} + 8\pi G\, T^{\Xi}_{\mu\nu}$.

\subsection{Vacuum/ground state}\label{sec:vacrenorm}

The potential $V(\Xi) = \Lambda^{4}\,\Xi\log\Xi$ has
\[
  V'(\Xi) = \Lambda^{4}(1 + \log\Xi),
  \qquad
  V''(\Xi) = \Lambda^{4}/\Xi,
\]
so $V'(\Xi_0) = 0$ at
\begin{equation}\label{eq:vacuum}
  \boxed{\;\Xi_0 = e^{-1} \approx 0.36788\;}
\end{equation}
and $V''(\Xi_0) = \Lambda^{4} e > 0$: $\Xi_0$ is a local minimum. The
location of this minimum at a coupling-independent dimensionless number
is generic for any potential of the form $\Lambda^{4} f(\Xi)$ with $\Xi$
dimensionless, since $V'(\Xi) = \Lambda^{4} f'(\Xi)$ vanishes
independently of $\Lambda$. What is specific to the choice
$f(\Xi) = \Xi\log\Xi$ is the value $\Xi_0 = e^{-1}$.

The potential energy at the vacuum is
\[
  V(\Xi_0) = \Lambda^{4}\,e^{-1} \log(e^{-1}) = -\Lambda^{4}/e < 0.
\]
This is the bare contribution from the scalar sector. As in standard
quintessence models, this combines with the bare cosmological constant
$\Lambda_{bare}$ from the gravitational sector to give the renormalized
observed cosmological constant $\Lambda_{obs} = \Lambda_{bare} - \Lambda^{4}/e$.
The bare quantities are not separately observable; only their sum enters
the Einstein equations. The scalar sector provides a dynamical
dark-energy component with a logarithmic potential and a vacuum
contribution whose observable effect is absorbed into the renormalized
cosmological constant. The model's empirical predictions for $w(z)$
depend on the dynamical evolution of $\Xi$ from $\Xi_i$ at $z_i$ toward
$\Xi_0 = e^{-1}$, not on the absolute value of $\Lambda_{bare}$.

\subsection{Fluctuation mass and stability}

Write $\Xi = \Xi_0 + \delta\Xi$ and linearize \eqref{eq:EL} about $\Xi_0$:
\[
  M_\Pl^{2}\,\Box\,\delta\Xi
  = \Lambda^{4}\,\frac{1}{\Xi_0}\,\delta\Xi
  = \Lambda^{4}\, e\,\delta\Xi.
\]
On a Minkowski background, a Fourier mode $\delta\Xi \propto e^{i(\mathbf{k}\cdot\mathbf{x} - \omega t)}$
satisfies $\omega^{2} - k^{2} = \Lambda^{4} e / M_\Pl^{2}$, yielding the
Klein-Gordon dispersion $\omega^{2} = k^{2} + m_\Xi^{2}$ with
\begin{equation}\label{eq:mass}
  \boxed{\;m_\Xi^{2} = \frac{\Lambda^{4}\,e}{M_\Pl^{2}}\;}
\end{equation}
Stability: $m_\Xi^{2} > 0$ for any $\Lambda > 0$.

\paragraph{Dark-energy scale.}
For $\Xi$ to play the role of dark energy with mass of order the Hubble
scale today, $m_\Xi \sim H_0 \sim 10^{-33}$ eV, equation \eqref{eq:mass}
implies $\Lambda \sim (H_0\, M_\Pl/e^{1/2})^{1/2}$, which evaluates to
$\Lambda \approx 1.5$~meV from this dimensional analysis. The numerical
fit reported in \S\ref{sec:desifit} gives $\Lambda \approx 1.7$~meV,
which differs from the dimensional estimate by a factor consistent with
the geometric factor relating $m_\Xi^{2}$ to $\rho_{c,0} = 3 H_0^{2}
M_\Pl^{2}$. Both values place $\Lambda$ numerically near the observed
dark-energy scale; this is a consistency observation rather than a
derivation.

\subsection{Global behavior of $V$}

For $\Xi > 0$, $V(\Xi) = \Lambda^{4}\,\Xi\log\Xi$ satisfies
\[
  V(\Xi) \to 0\ \text{from below as}\ \Xi \to 0^{+},\qquad
  V(\Xi_0) = -\Lambda^{4}/e\ (\text{global minimum}),\qquad
  V(\Xi) \to +\infty\ \text{as}\ \Xi \to +\infty.
\]
The scalar potential is globally bounded below on $\Xi > 0$, and the
homogeneous scalar-sector energy density
$\rho_\Xi = \tfrac{M_\Pl^{2}}{2}\dot\Xi^{2} + V(\Xi)$ is bounded below
by $-\Lambda^{4}/e$.

\section{FRW Cosmology}\label{sec:frw}

\subsection{Background equations}

In a spatially flat FRW background with metric
$ds^{2} = -dt^{2} + a^{2}(t)\, d\mathbf{x}^{2}$ and homogeneous $\Xi(t)$,
the kinetic invariant is
$g^{\alpha\beta}\partial_\alpha \Xi\, \partial_\beta \Xi = -\dot\Xi^{2}$.
Writing $T^{\Xi}_{\mu\nu} = \diag(\rho_\Xi, a^{2} p_\Xi, a^{2} p_\Xi, a^{2} p_\Xi)$,
\begin{align}
  \rho_\Xi &= \frac{M_\Pl^{2}}{2}\,\dot\Xi^{2} + \Lambda^{4}\,\Xi\log\Xi,\label{eq:rho}\\
  p_\Xi &= \frac{M_\Pl^{2}}{2}\,\dot\Xi^{2} - \Lambda^{4}\,\Xi\log\Xi.\label{eq:p}
\end{align}
The Friedmann equations are
\begin{equation}\label{eq:friedmann}
  H^{2} = \frac{8\pi G}{3}\!\left(\rho_{\SM} + \rho_\Xi\right),\qquad
  \dot H = -4\pi G\!\left(\rho_{\SM} + p_{\SM} + M_\Pl^{2}\,\dot\Xi^{2}\right),
\end{equation}
and the $\Xi$ equation of motion, from \eqref{eq:EL} in FRW (where
$\Box\Xi = -\ddot\Xi - 3H\dot\Xi$ in the $(-,+,+,+)$ convention), is
\begin{equation}\label{eq:xieom}
  \boxed{\;M_\Pl^{2}(\ddot\Xi + 3H\dot\Xi) = -\Lambda^{4}(1 + \log\Xi)\;}
\end{equation}

\subsection{Equation of state}

From \eqref{eq:rho}-\eqref{eq:p},
\begin{equation}\label{eq:w}
  w_\Xi = \frac{p_\Xi}{\rho_\Xi}
  = \frac{\tfrac{1}{2} M_\Pl^{2}\dot\Xi^{2} - \Lambda^{4}\,\Xi\log\Xi}
         {\tfrac{1}{2} M_\Pl^{2}\dot\Xi^{2} + \Lambda^{4}\,\Xi\log\Xi}
  = -1 + \frac{M_\Pl^{2}\dot\Xi^{2}}{\rho_\Xi}.
\end{equation}

\begin{proposition}[Limiting regimes]\label{prop:limits}
For $\Xi > 0$ the EoS \eqref{eq:w} satisfies
\begin{enumerate}
  \item[(i)] $w_\Xi = -1$ at the vacuum ($\dot\Xi = 0$, $\Xi = \Xi_0$):
        vacuum-like equation of state.
  \item[(ii)] On the physical rolling branch studied numerically
        (\S\ref{sec:desifit}), $w_\Xi$ approaches $-1$ from above as
        $\Xi \to \Xi_0$ and $\dot\Xi \to 0$.
  \item[(iii)] $w_\Xi \to +1$ in the kinetic-domination limit
        ($M_\Pl^{2}\dot\Xi^{2} \gg \Lambda^{4}|\Xi\log\Xi|$): stiff fluid.
\end{enumerate}
\end{proposition}

\begin{proof}
(i) Direct substitution in \eqref{eq:w}: $p_\Xi = -\Lambda^{4}/e$ and
$\rho_\Xi = \Lambda^{4}\,\Xi_0 \log\Xi_0 = -\Lambda^{4}/e$, giving
$w_\Xi = -1$. (ii) On the monotone rolling solutions studied numerically,
with $\Xi > \Xi_0$ and $\rho_\Xi > 0$ over the fitted redshift range, the
kinetic term decreases as the field approaches the minimum, and
\eqref{eq:w} gives $w_\Xi \to -1$ from above.
(iii) In the kinetic-domination regime, both numerator and denominator
of \eqref{eq:w} are dominated by $M_\Pl^{2}\dot\Xi^{2}$, giving
$w_\Xi \to +1$.
\end{proof}

\begin{remark}[Bare vs renormalized vacuum]
The equality $w_\Xi = -1$ at $\Xi = \Xi_0$ in (i) is a formal statement
about the scalar-sector stress-energy tensor before vacuum-energy
renormalization; the physical late-time vacuum entering cosmology is the
renormalized sum discussed in \S\ref{sec:vacrenorm}.
\end{remark}

\begin{remark}[On phantom crossing]
A general structural theorem $w_\Xi \geq -1$ would require careful
treatment of the bare scalar-sector energy density, which is negative
at the vacuum ($\rho_\Xi(\Xi_0) = -\Lambda^4/e$). Over the fitted
rolling branch in \S\ref{sec:desifit}, the numerical solution exhibits
$w_\Xi(z) \geq -1$ at all redshifts and produces no phantom crossing;
we treat the absence of phantom crossing as a property of the fitted
branch rather than a global theorem of the minimal theory.
Phantom-crossing dark energy ($w < -1$) and its consequences are
treated in a separate literature \cite{Caldwell2002Phantom}, which is
not the regime of the present model on the rolling branch.
\end{remark}

\subsection{Late-time attractor}

Combining \eqref{eq:xieom} with \eqref{eq:friedmann}, the field rolls
toward the vacuum, and the damping term $3H\dot\Xi$ drives $\dot\Xi \to 0$
as $H$ remains positive. In this limit $\Xi \to \Xi_0 = e^{-1}$ and, by
Proposition \ref{prop:limits}(i), $w_\Xi \to -1$. The
late-time attractor is therefore vacuum-like behavior with $w \to -1$.
This belongs to the freezing class in the standard quintessence
classification \cite{CaldwellLinder2005,ScherrerSen2008Thawing}; the
distinguishing feature of the present model is not the late-time
attractor itself but the dual-regime traversal en route to it
(\S\ref{sec:typeTFA}).

\section{Structural Connections to Information-Theoretic Functionals}\label{sec:entropy}

The functional form $\Xi\log\Xi$ that appears in the potential
\eqref{eq:Vxi} is recognizable from two adjacent settings: the per-bin
integrand of the Gibbs/Shannon entropy on a normalized probability
distribution, and the Bialynicki-Birula--Mycielski nonlinearity in
nonlinear quantum mechanics. We treat both connections as structural
observations about the functional form of $V$, not as derivations of
\eqref{eq:action}, and we are explicit below about where the analogies
break down.

\subsection{The Gibbs functional}

The function
\begin{equation}\label{eq:entropy}
  H_{\Gibbs}(x) = -x \log x,
  \qquad x \in (0,1],
\end{equation}
is the per-bin contribution to the Gibbs/Shannon entropy of a discrete
probability distribution $\{p_i\}$ via $H = \sum_i H_{\Gibbs}(p_i)$,
where the $p_i$ are nonnegative and sum to $1$. The function attains its
maximum on $(0,1]$ at $x = e^{-1}$, with $H'_{\Gibbs}(e^{-1}) = 0$ and
$H''_{\Gibbs}(e^{-1}) = -e < 0$.

The potential \eqref{eq:Vxi} has the same functional form up to a sign
and the dimensionful prefactor:
\begin{equation}\label{eq:formal-rel}
  V(\Xi) = -\Lambda^{4}\,H_{\Gibbs}(\Xi).
\end{equation}
The vacuum condition $V'(\Xi_0) = 0$ is therefore equivalent to
$H'_{\Gibbs}(\Xi_0) = 0$, and the same $\Xi_0 = e^{-1}$ that minimizes
$V$ also maximizes $H_{\Gibbs}$ on $(0,1]$.

\begin{remark}[Where the analogy breaks]
The function $H_{\Gibbs}$ as written has an operational interpretation
as a thermodynamic entropy only when $x$ is a normalized probability
($x \in [0,1]$ with $\sum x = 1$). In the present scalar-field theory
$\Xi$ is a real, positive, unbounded scalar with no normalization
constraint. The relation \eqref{eq:formal-rel} is therefore a
structural rhyme between $V$ and the Gibbs integrand, not an
identification of $-V/\Lambda^{4}$ with a thermodynamic entropy.
We do not claim that the field's dynamics is equivalent to entropy
maximization in any operational sense; we record only that the
potential's functional form coincides with the Gibbs integrand and
that $\Xi_0 = e^{-1}$ coincides numerically with the per-bin
entropy maximum. The model's empirical predictions for $w(z)$
\S\ref{sec:obs} do not depend on this connection.
\end{remark}

\subsection{Bialynicki-Birula--Mycielski nonlinearity}

The same logarithmic functional form appears in nonlinear quantum
mechanics: Bialynicki-Birula and Mycielski
\cite{BialynickiBirulaMycielski1976} characterized
$|\psi|^{2} \log |\psi|^{2}$ as the unique analytic nonlinear extension
of the free Schr\"odinger equation that preserves separability of the
density $|\psi|^{2}$ under wave-function factorization
$\psi = \psi_1 \psi_2$. The same logarithmic functional appears in
nuclear physics applications of nonlinear Schr\"odinger dynamics
\cite{Hefter1985}, and precision laboratory tests of nonlinear quantum
mechanics with trapped ions and atom interferometry have constrained
candidate nonlinearities of this and related forms
\cite{Weinberg1989Precision,Bollinger1989}.

\begin{remark}[Where the analogy breaks]
The Bialynicki-Birula--Mycielski nonlinearity is in
$|\psi|^{2}\log|\psi|^{2}$, where $|\psi|^{2}$ is a probability density.
In the present setting, the analogous variable is the field $\Xi$
itself, not $\Xi^{2}$, and $\Xi$ has no probability-density
interpretation. Importing the same functional form from a setting
where the argument is a normalized squared amplitude to a setting where
the argument is an unbounded real scalar is a structural rhyme, not a
derivation of \eqref{eq:action} from the BBM uniqueness theorem. We
cite \cite{BialynickiBirulaMycielski1976,Rosen1969,CazenaveHaraux1980,HoeghKrohn1971}
as motivation for studying $V \propto \Xi\log\Xi$ as a candidate
quintessence potential, not as a uniqueness proof for it.
\end{remark}

\subsection{Logarithmic nonlinearity in cosmology: a related research lineage}\label{sec:loglineage}

A substantial body of prior work has explored the logarithmic
Schr\"odinger / Klein-Gordon equation in cosmological and dark-sector
contexts. Zloshchastiev's program treats
$|\psi|^{2}\log|\psi|^{2}$-type nonlinearity as a fundamental
correction to quantum field theory with cosmological consequences
\cite{Zloshchastiev2010LogTime,Zloshchastiev2011SSB,AvdeenkovZloshchastiev2011,Zloshchastiev2018Stability,Zloshchastiev2020SuperfluidDM},
with applications including emergence of evolution time from a
generally-covariant logarithmic theory
\cite{Zloshchastiev2010LogTime}, spontaneous symmetry breaking in
logarithmic nonlinear quantum theory \cite{Zloshchastiev2011SSB},
self-sustainability and emergence of spatial extent in quantum Bose
liquids with logarithmic nonlinearity
\cite{AvdeenkovZloshchastiev2011}, and superfluid-vacuum-theory
treatments of dark matter and dark energy
\cite{Zloshchastiev2020SuperfluidDM}. Chavanis's logotropic-dark-fluid
program treats logarithmic equations of state as candidate unifications
of dark matter and dark energy
\cite{Chavanis2015,Chavanis2021Kessence,Chavanis2022Logotropic}, with
the most recent installment \cite{Chavanis2022Logotropic} introducing a
\emph{complex} scalar field with a logarithmic potential satisfying a
generalized Klein-Gordon equation, returning $\Lambda$CDM in the
classical limit $\hbar \to 0$, and unifying the dark sector by
identifying its rest-mass content as dark matter and its internal
energy as dark energy.

The present work shares the use of a logarithmic potential with these
programs but differs from each in three respects relevant to the
referee. (i) \emph{Field content.} The present paper studies a real,
positive, dimensionless scalar $\Xi$ with no normalization or
density-amplitude interpretation; Zloshchastiev's program is in
$|\psi|^{2}$ where $|\psi|^{2}$ is a probability density and the
logarithmic nonlinearity sits inside the wave equation rather than as
a potential term in a classical scalar action; Chavanis's recent
\cite{Chavanis2022Logotropic} construction uses a complex scalar field
$\varphi$ with potential $V(|\varphi|^{2})$ involving a logarithm of
$|\varphi|^{2}$, with an explicit $\hbar \to 0$ classical limit
returning $\Lambda$CDM. (ii) \emph{Sector unification.} The present
paper proposes $\Xi$ as a dark-energy field only and makes no claim of
dark-matter unification; the logotropic and superfluid-vacuum programs
explicitly aim at DM/DE unification. (iii) \emph{Dynamical signature.}
The present paper's central observational claim is the dual-regime
freeze-thaw transit (\S\ref{sec:typeTFA}) with a Type-F frozen
turnaround at intermediate redshift $z_\star$ and the resulting
non-monotone $w_{\rm DE}(z)$ signature
(\S\ref{sec:falsify}, criterion (F$_5$)); to our knowledge this
trajectory class has not been singled out as the central
phenomenological prediction in the prior logarithmic-cosmology
literature, which has largely focused on $\Lambda$CDM-equivalence at
large scales and structure-formation phenomenology at galactic scales.

We cite the prior logarithmic-nonlinearity-in-cosmology lineage as
substantive related work, not as derivational input to the present
action. Nothing in the empirical content of \S\ref{sec:obs} depends
on the structural relations to that lineage.

\subsection{Information-theoretic precedents}

\begin{remark}[Information-theoretic dark energy]
Entropic approaches to dark energy in the holographic sense
\cite{Verlinde2011,CHLZ2012} typically start from an entropy ansatz
and derive a dark-energy density. The present approach is structurally
different: the action \eqref{eq:action} is specified first, and the
relation \eqref{eq:formal-rel} between $V$ and the Gibbs integrand
follows. Information-theoretic field theory \cite{Ensslin2013,Caticha2012}
uses $\Xi\log\Xi$ as a standard KL-divergence integrand on spaces of
probability densities; the appearance of the same functional as a
cosmological self-interaction is a point of contact between these
literatures and the present work, but the connection is again
structural rather than derivational.
\end{remark}

The model's empirical content (\S\ref{sec:obs}) does not depend on either
of the two structural connections discussed in this section.

\section{Observational Predictions}\label{sec:obs}

\subsection{The $w(z)$ trajectory}

On the physical rolling branch studied numerically, we write
$w_\Xi(z) = -1 + \epsilon(z)$ with $\epsilon(z) \geq 0$. The freezing
attractor derived above gives $\epsilon(z) \to 0$ at late times, with
the rate determined by $\Lambda$ and the initial displacement $\Xi(z_i)$
at some reference redshift $z_i$. DESI DR1/DR2 measurements of $w_0$
and $w_a$ in the CPL parametrization probe $\epsilon(z)$ directly at
$z \lesssim 1.5$ \cite{DESI2024BAO,DESI2024VI}.

\subsection{Numerical integration and DESI consistency check}\label{sec:desifit}

We perform an illustrative background-level consistency check against
the DESI~2024 DR1 published $(w_0, w_a)$ central values, treating those
values as a Gaussian summary on the CPL parametrization. We emphasize at
the outset what this check does and does not do: the
$\chi^{2}$ values reported below are computed against the published
$(w_0, w_a)$ marginal Gaussian summary, with central values
$w_0 = -0.827 \pm 0.063$, $w_a = -0.75 \pm 0.27$ \cite{DESI2024VI}, and
they quantify how close the model's $(w_0, w_a)$ output sits to those
published central values \emph{under that Gaussian-on-summary
approximation}. They are not derived from the underlying joint
BAO + CMB + SN likelihood. The DR1 BAO data alone are consistent with
flat $\Lambda$CDM \cite{DESI2024BAO}; the dynamical-dark-energy
preference at the $2.8$--$4.2\sigma$ level appears in the joint
likelihood, not in the marginal Gaussian on $(w_0, w_a)$. A full joint
likelihood analysis using the DR2 chains \cite{DESI2025DR2} together
with Pantheon+ supernovae and Planck CMB constraints is in preparation
as a companion numerical paper.

We numerically integrate the FRW system
\eqref{eq:friedmann}-\eqref{eq:xieom} over $z \in [0, 10]$ with a
standard-density background ($\Omega_m = 0.315$, $\Omega_r = 9.1\times10^{-5}$,
$\Omega_\Xi = 0.6849$, $H_0 = 67.4$ km s$^{-1}$ Mpc$^{-1}$
\cite{Planck2020}) and scan over $(\Lambda^{4}/\rho_{c,0}, \Xi_i, \dot\Xi_i)$
specified at the reference redshift $z_i \approx 20$, where $\rho_{c,0} = 3 H_0^{2} M_\Pl^{2}$
is the critical density today.\footnote{For numerical stability, the supplementary
script \texttt{desi\_xi\_optimize\_v2.py} integrates from an earlier matter-era
starting point at $N_\mathrm{start} = -4$ (corresponding to $z \approx 54$), then
reads off the trajectory state at $z = 20$ to define the documented initial-condition
values $(\Xi_i, \dot\Xi_i)$. Equivalently, one can integrate forward from
$z_i = 20$ using cosmic-time variables; both prescriptions yield the same
late-time evolution of the rolling branch. The reconciliation script
\texttt{compute\_zstar\_v3.py} demonstrates the equivalence by reproducing the
trajectory in either convention.}
A notational convention used throughout: the script reports densities in
$\Omega$-units (each density pre-divided by $\rho_{c,0} = 3 H_0^{2} M_\Pl^{2}$),
so the implicit Friedmann equation in the script reads $H^{2}(1 - \tfrac{1}{2}\Xi'^{2}) =
\Omega_m a^{-3} + \Omega_r a^{-4} + \Lambda^{4}\Xi\log\Xi$ where $\Xi' = d\Xi/dN$.
This is algebraically equivalent to the standard form $3 H^{2} = \rho_m + \rho_r +
\tfrac{1}{2}M_\Pl^{2}\dot\Xi^{2} + V(\Xi)$ with $\dot\Xi = H\Xi'$ and physical
densities; the factor of 3 is absorbed into the $\Omega$-unit conversion. We restrict to the rolling branch of the
model on which $w_\Xi(z) \geq -1$ throughout the integrated range
(criterion (F$_1$) of \S\ref{sec:falsify}) and on which the present-epoch
inbound leg is monotone over $0 \leq z \leq 1.5$ (criterion (F$_2$);
\S\ref{sec:typeTFA} for the full Type-T/F/A trajectory taxonomy).
A configuration that yields proximity to the DESI DR1 published
$(w_0, w_a)$ Gaussian summary is
\begin{equation}\label{eq:desibestfit}
  \Lambda^{4}/\rho_{c,0} = 0.231,
  \quad \Xi_i = 0.925,
  \quad \dot\Xi_i = +0.429,
\end{equation}
corresponding to $\Lambda \approx 1.7\,\text{meV}$ in physical units and
yielding
\[
  w_0^{\Xi} = -0.793,\quad
  w_a^{\Xi} = -0.451,\quad
  \chi^{2}_{\text{Gauss}} = 1.52,
\]
to be compared with $\chi^{2}_{\text{Gauss},\Lambda CDM} = 15.26$ for
$(w_0, w_a) = (-1, 0)$ under the same Gaussian-on-summary
approximation. We use the subscript ``Gauss'' to mark these
$\chi^{2}$ values as quantifying proximity to the published
$(w_0, w_a)$ marginal Gaussian, not goodness of fit to the underlying
data. The $\Xi$-model's $w(z)$ trajectory on this configuration is
monotonically inbound on $0 \leq z \leq z_\star$ with $z_\star \approx 2$
the Type-F turnaround (cf.\ \S\ref{sec:typeTFA}):

\begin{center}
\begin{tabular}{c|c}
$z$ & $w_\Xi(z)$\\
\hline
$0.0$ & $-0.7933$\\
$0.3$ & $-0.9098$\\
$0.5$ & $-0.9432$\\
$0.8$ & $-0.9701$\\
$1.0$ & $-0.9804$\\
$1.5$ & $-0.9946$\\
$2.0$ & $-0.9998$
\end{tabular}
\end{center}

The present-epoch field value is $\Xi(z=0) \approx 0.97$, still above
the vacuum $\Xi_0 = e^{-1} \approx 0.37$, confirming the field is in
its rolling phase. The numerical integration is performed by the
script \texttt{desi\_xi\_optimize\_v2.py} provided as electronic
supplementary material.

\paragraph{Initial-condition tuning.}
The values $(\Xi_i, \dot\Xi_i) = (0.925, +0.429)$ at $z_i \approx 20$ in
\eqref{eq:desibestfit} are tuned, not predicted. They are selected by
the requirement of proximity to the DESI DR1 $(w_0, w_a)$ central values
together with the rolling-branch criteria $w_\Xi(z) \geq -1$ and
monotone inbound $w_\Xi(z)$ on $0 \leq z \leq 1.5$ ((F$_1$), (F$_2$) of
\S\ref{sec:falsify}). We note explicitly that $\dot\Xi_i > 0$
corresponds to an \emph{outbound} initial condition: at $z_i$ the
field is moving away from the vacuum $\Xi_0 = e^{-1}$, decelerates
under Hubble friction and potential gradient to a Type-F turnaround at
$z_\star \approx 2$ (\S\ref{sec:typeTFA}), and rolls back toward $\Xi_0$
on the inbound leg containing the present epoch. The direction of
initial motion is itself a tuning choice in addition to the values
$(\Xi_i, \dot\Xi_i)$. There is no naturalness or attractor argument in
the present minimal theory that singles out these values; the
rolling-branch restriction selects solutions, it does not derive them.
This sensitivity to initial conditions is shared with most
quintessence models that fit late-time data, and a complete
naturalness analysis (e.g., a quintessence-tracker construction) is
outside the scope of the present paper. A reduction of this
two-dimensional initial-condition freedom to a one-dimensional or
zero-dimensional family --- e.g., via a late-time attractor mechanism
that erases initial-condition memory --- would strengthen the model
and is a target for follow-up work.

\paragraph{Methodological scope of this consistency check.}
To restate the scope explicitly: the $\chi^{2}_{\text{Gauss}}$ values
reported here are computed against published CPL summary statistics,
treated as an uncorrelated Gaussian likelihood on $(w_0, w_a)$. The
DESI DR1/DR2 publications provide non-negligible $(w_0, w_a)$
correlation in their full chains; ignoring that correlation is one
limitation of the present check. A full joint likelihood analysis
on the underlying BAO + CMB + SN data is deferred to a companion
numerical paper. The present check demonstrates that the model can
produce $(w_0, w_a)$ values within $\sim 1\sigma$ of the DESI DR1
central values for $\Lambda \approx 1.7$~meV with appropriately
chosen initial conditions, and that $\Lambda$CDM sits roughly
$3.9\sigma$ from those central values under the same Gaussian
approximation. It is not a claim of best-fit to the full DESI
likelihood, nor a measure of the model's preference relative to
$\Lambda$CDM in any joint-likelihood sense.

\paragraph{Reproducibility.}
The numerical fit reported above is reproduced by a single command. The
script \texttt{desi\_xi\_optimize\_v2.py} (electronic supplementary material)
accepts no arguments; running it produces the best-fit parameters
\eqref{eq:desibestfit}, the $w(z)$ trajectory, and the
$\chi^{2}_{\text{Gauss}}$ values for both the $\Xi$-model and
$\Lambda$CDM. Total runtime is under 30 seconds on a standard laptop.
The verification script \texttt{proof\_xi\_canonical.py} runs 22
algebraic and stability tests on the field equations; all pass. Both
scripts are archived at DOI~\texttt{10.5281/zenodo.18852047}.

\subsection{Type-T/F/A classification of the trajectory}\label{sec:typeTFA}

The cosmological trajectory of the field $\Xi$ in the documented fit
\eqref{eq:desibestfit} traverses three dynamically distinct regimes
within a single physical history. We classify these as follows.

\paragraph{Type-T (Thawing regime).}
The field is displaced from the vacuum ($\Xi$ noticeably different from
$\Xi_0 = e^{-1}$) with non-negligible kinetic energy, $\dot\Xi \neq 0$.
The bare scalar EoS $w_\Xi(z)$ deviates from $-1$ by an amount controlled
by the ratio of the scalar's kinetic energy to its full energy density
\eqref{eq:w}, and the combined-sector effective $w_{\rm DE}(z)$ inherits
this deviation through the standard quintessence formula. In the
documented fit, the current cosmological epoch $z = 0$ is in the Type-T
regime: the field is on the inbound leg of the trajectory, having
passed through the turnaround near $z \approx 2$, and
$w_\Xi(z = 0) \approx -0.79$ reflects the kinetic-to-total ratio of
the scalar sector at the present epoch.

\paragraph{Type-F (Frozen turnaround).}
At a specific redshift $z_\star$ on the trajectory, the field reaches
its outbound maximum and $\dot\Xi = 0$ momentarily, with the bare
scalar EoS $w_\Xi(z_\star) = -1$ instantaneously by direct evaluation
of \eqref{eq:w}. The combined-sector $w_{\rm DE}(z_\star)$ takes its
local minimum value in the neighborhood of $z_\star$, approaching
$-1$ from above. In the documented fit this minimum sits at
$z_\star \approx 2$, with $w_\Xi(z_\star) \approx -0.9998$ in the
tabulated values of \S\ref{sec:desifit}; the residual $\sim 10^{-4}$
deviation from $-1$ in the tabulated value reflects finite numerical
sampling near the turnaround rather than physical departure from the
analytic Type-F condition $\dot\Xi(z_\star) = 0$. The Type-F state is
not a fixed point of the equation of motion: the field has zero
velocity but nonzero potential gradient ($V'(\Xi_{\max}) > 0$ since
$\Xi_{\max} > \Xi_0$), and the next instant the field begins to roll
inbound under the potential force, modulated by Hubble friction.

\paragraph{Type-A (Asymptotic refreeze).}
In the cosmological future ($z \to -1$, $a \to \infty$), the field
rolls back toward $\Xi_0 = e^{-1}$ with Hubble friction damping the
inbound motion. As $\Xi \to \Xi_0$ and $\dot\Xi \to 0$, both $\rho_\Xi$
and $V_{\rm dev}(\Xi) := V(\Xi) - V(\Xi_0)$ vanish, and $w_\Xi(z) \to -1$
from above. This is the late-time attractor identified in
\S\ref{sec:frw}: the universe approaches an effective $\Lambda$CDM
state with cosmological constant set by the renormalized vacuum
energy.

\paragraph{The geometric anchor at $\Xi_0 = e^{-1}$.}
The vacuum value $\Xi_0 = e^{-1}$ plays three independent roles in the
model: (a) the unique global minimum of the potential
$V(\Xi) = \Lambda^4\,\Xi\log\Xi$, (b) the unique maximizer of the
per-bin Gibbs functional $H_{\rm Gibbs}(x) = -x\log x$ on $(0, 1]$, and
(c) the geometric anchor of the freeze-thaw transit --- the Type-A
endpoint is exactly $\Xi_0$, and the Type-F turnaround at $\Xi_{\max}$
is the locus where the integrated outbound kinetic energy is exactly
balanced by the work done against $V'(\Xi)$ along the outbound leg
referenced to $\Xi_0$. The three characterizations are independent:
(a) is an analytic property of $V$, (b) is a structural correspondence
with the per-bin Gibbs integrand, and (c) is a dynamical property of
the FRW trajectory on the rolling branch. We do not claim that the
same numerical value $e^{-1}$ playing these three roles constitutes a
derivation of the action --- only that the same geometric anchor
organizes the potential's analytic structure, its information-theoretic
content, and its cosmological dynamics simultaneously.

\paragraph{Distinction from standard freezing/thawing quintessence.}
Caldwell-Linder \cite{CaldwellLinder2005} classify quintessence
trajectories into two mutually exclusive dynamical regimes:
\emph{freezing}, with monotone $w(z) \to -1$ from above as $z$
decreases and no local minimum, and \emph{thawing}, with monotone
departure from $-1$ over the observable epoch and no return. The
dual-regime trajectory studied here belongs to neither class: it
traverses Type-T (post-turnaround inbound thaw on the present epoch),
Type-F (instantaneous frozen turnaround at $z_\star \approx 2$), and
Type-A (asymptotic refreeze in the cosmological future) within a
single physical history. The observational signature distinguishing
this dual-regime trajectory from single-regime quintessence is the
existence of a local minimum in $w_{\rm DE}(z)$ at intermediate
redshift, with the minimum value approaching $-1$ from above.
Falsification criterion $(F_5)$ in \S\ref{sec:falsify} makes this
prediction explicit and testable at Stage-IV precision.

\paragraph{Sensitivity of $z_\star$ to fit parameters.}
The location of the Type-F turnaround $z_\star$ depends on the
trajectory through $(\Lambda^4, \Xi_i, \dot\Xi_i)$. The documented fit
\eqref{eq:desibestfit} gives $z_\star \approx 2$. A sensitivity scan
on the optimization script's local neighborhood gives
$1.5 \lesssim z_\star \lesssim 3.0$ over the manuscript's $1\sigma$
parameter region; the precise sensitivity range is reported by the
optimization script
\texttt{desi\_xi\_optimize\_v2.py} (electronic supplementary material).
This range is the support over which $(F_5)$ must be tested at
Stage-IV precision; trajectory paths with $z_\star$ substantially
outside this range require correspondingly different
$(\Lambda^4, \Xi_i, \dot\Xi_i)$ that are excluded by the
$(w_0, w_a)$ proximity check.

\subsection{Falsification criteria}\label{sec:falsify}

The model predicts:
\begin{enumerate}
  \item[(F1)] On the monotone rolling branch selected by the present
        model, $w_\Xi(z) \geq -1$ over the redshift range studied
        numerically. A robust observational preference for $w(z) < -1$
        across that same range, at $\geq 3\sigma$ confidence in the joint
        DR2 + Pantheon+ + Planck likelihood, would disfavor this branch
        of the model.
  \item[(F2)] Monotone inbound approach $w_\Xi(z) \to -1$ from above
        on the redshift interval below the Type-F turnaround
        ($0 \leq z \lesssim z_\star$, with $z_\star \approx 2$ in the
        documented fit). The current cosmological epoch sits on the
        inbound leg of the trajectory; on this leg the
        $w(z)$ profile is monotone in $z$. A confirmed non-monotone
        $w(z)$ on $0 \leq z \leq 1.5$ at $\geq 3\sigma$ in the joint
        DR2 + Pantheon+ + Planck likelihood would falsify the
        rolling-branch fit on the inbound leg. (Note: the global
        trajectory is non-monotone through the Type-F turnaround at
        $z_\star$; (F2) constrains only the inbound interval below
        $z_\star$. The global non-monotone signature is criterion
        $(F_5)$ below.)
  \item[(F3)] Asymptotic vacuum endpoint $w \to -1$. The present-epoch
        value $w_0 \neq -1$ in general while the field is still rolling;
        the manuscript's best fit gives $w_0 \approx -0.79$ at the present
        epoch. Evidence favoring late-time behavior inconsistent with
        $w \to -1$, at $\geq 3\sigma$ in the joint Stage-IV likelihood,
        would disfavor the freezing attractor interpretation.
  \item[(F4)] Specific $w(z)$ profile fixed by two parameters
        $(\Lambda^{4}, \Xi_i)$. The model predicts a structured
        trajectory $w_\Xi(z)$ that, after fitting $w_0$ and $w_a$, is
        overdetermined: any third independent CPL-equivalent moment
        (e.g., $dw/dz$ at $z=0.5$, or $w$ at $z=1.0$) is a non-trivial
        test.
  \item[(F5)] \emph{Non-monotone $w_{\rm DE}(z)$ with local minimum
        near $-1$ at intermediate redshift.}
        The dual-regime transit predicts that the combined-sector
        effective $w_{\rm DE}(z)$ reaches a local minimum at
        $z = z_\star$ corresponding to the Type-F turnaround
        (\S\ref{sec:typeTFA}), with $w_{\rm DE}(z_\star) \to -1$ from
        above. In the documented fit \eqref{eq:desibestfit},
        $z_\star \approx 2$, with the precise value determined by
        $(\Lambda^{4}, \Xi_i, \dot\Xi_i)$ and shifting within
        $1.5 \lesssim z_\star \lesssim 3.0$ over the manuscript's
        $1\sigma$ parameter region. Standard freezing quintessence
        \cite{CaldwellLinder2005} predicts monotone $w(z) \to -1$ from
        above with no local minimum; standard thawing quintessence
        predicts monotone departure from $-1$ with no return. Only a
        trajectory that traverses both regimes admits a local minimum.
        \emph{Falsification:} a reconstructed $w_{\rm DE}(z)$ that is
        monotonic over $0 \leq z \leq 5$ at $\geq 3\sigma$ confidence
        in the joint DR2 + Pantheon+ + Planck likelihood, with no
        local minimum detected within the Stage-IV measurement
        uncertainty $\sigma(w(z)) \approx 0.03$
        \cite{ShajibFrieman2025}, would disfavor the dual-regime
        transit interpretation, leaving (within this model) only the
        non-rolling-branch configurations that do not realize the
        Type-F turnaround on the observable interval.
\end{enumerate}
Stage-IV surveys (DESI through DR3, Euclid \cite{EuclidLaureijs2011},
Rubin LSST \cite{LSSTIvezic2019}) will reduce
$\sigma(w_0)$ to approximately $0.02$ and $\sigma(w_a)$ to approximately
$0.10$ \cite{ShajibFrieman2025,Turyshev2026Review}, providing the
precision needed to test (F1)--(F5) at the levels indicated above.
Forecasting tools and Stage-IV reach are reviewed in
\cite{Albrecht2006DETF,Aubourg2015}.
The dual-regime signature (F5) requires non-parametric $w(z)$
reconstructions beyond the CPL parametrization, since CPL is monotone
by construction and cannot represent a local minimum at
$z_\star \approx 2$. Standard non-parametric techniques applicable
to the F5 test include principal-component analysis of binned $w(z)$
\cite{CrittendenPogosianZhao2009}, Gaussian-process reconstruction
\cite{Holsclaw2010GP}, the $Om$ and related null diagnostics
\cite{SahniShafielooStarobinsky2008}, and tomographic reconstruction
\cite{PogosianShafieloo2009}. Alternative parametrizations beyond CPL
that admit a single inflection point or finite-time turnaround
\cite{JassalBaglaPadmanabhan2005,BarbozaAlcaniz2008} can capture the
Type-F signature at the parametric level, although they do not by
themselves predict its location. Recent DR2-era reconstructions
\cite{Wang2026QReconstruction,Adil2026Analytic} indicate the
precision and methodology relevant to this test; field-space
parametrizations of deviations from $\Lambda$
\cite{CrittendenMajerottoPiazza2007,Linder2006Paths} provide
complementary diagnostics tied to scalar-field dynamics.

\subsection{Fifth-force and local tests}

In \eqref{eq:action}, $\Xi$ has no direct matter coupling; any effective
coupling to SM matter arises only through gravity and is Planck-suppressed.
Such gravitational-strength couplings are well below the sensitivity of
laboratory tests of the inverse-square law \cite{EotWash,MICROSCOPE},
which probe direct-matter couplings several orders of magnitude above
the gravitational scale. The minimal theory is therefore expected to
produce \emph{no} observable fifth force and is consistent with all
current local tests. Any direct-matter coupling (conformal
$g\Xi\, T^{\SM}$, disformal $f(\Xi) T^{\SM}$, or universal rescaling
$(\Xi/\Xi_0) \mathcal{L}_{\SM}$) is outside the scope of the minimal
action and would be added as a separate extension.

\subsection{Effective-field-theory validity over the field
excursion}\label{sec:eft}

In canonical variables $\phi = M_\Pl \Xi$ (cf.\ the footnote in
\S\ref{sec:fieldcontent}), the vacuum sits at $\phi_0 = M_\Pl/e$ and
the fitted trajectory has $\phi$ traversing field-space distances of
order $M_\Pl$. This is the standard quintessence transplanckian-field
concern: an EFT obtained by integrating out heavier degrees of freedom
at a scale $M$ generically receives operator corrections suppressed by
powers of $\phi/M$, and for $\phi \sim M_\Pl$ such corrections are not
manifestly under control unless the underlying theory protects them
(via approximate shift symmetry, supersymmetry, or some other
mechanism). The minimal action \eqref{eq:action} does not exhibit a
shift symmetry, since $\Xi\log\Xi$ is not invariant under
$\Xi \to \Xi + c$, and we do not propose an underlying UV completion in
the present paper.

We make three observations. (i) The field excursion in canonical
variables is transplanckian in absolute terms: even within the
fitted late-time regime ($0 \le z \le 2$), $\Delta\phi$ is of order
$10^{-1}\,M_\Pl$, and the full evolution from $z_i \approx 20$
to today traverses $\Delta\phi \sim M_\Pl$. The argument that the
EFT remains controlled cannot rely on a small absolute field
excursion. We do not have a UV completion that protects the form
of $V$ against $M_\Pl$-scale operator corrections, and the minimal
action \eqref{eq:action} does not exhibit a shift symmetry that
would suppress them in the absence of such a completion. The
present paper treats \eqref{eq:action} as a tree-level effective
description; the question of UV control over the transplanckian
field excursion is open. This is the standard transplanckian
issue for any quintessence model with $\Delta\phi \sim M_\Pl$ and
is shared with most freezing-quintessence constructions in the
literature.
(ii) The mass scale $\Lambda \approx 1.7$~meV is many
orders of magnitude below $M_\Pl$, and the dimensional analysis
$m_\Xi^{2} = \Lambda^{4} e / M_\Pl^{2} \sim H_0^{2}$ does not require
new physics at scales below $M_\Pl$. (iii) A controlled UV embedding
of \eqref{eq:action} is a standard open question for any tree-level
quintessence potential and is not addressed here. We flag this as a
limitation and leave it to subsequent work; the present paper studies
the phenomenology of \eqref{eq:action} treated as a low-energy
effective description.

\subsection{Perturbations}\label{sec:perts}

The minimally coupled canonical scalar field \eqref{eq:action} has
\emph{rest-frame} sound speed
\[
  c_{s}^{2} \;=\; 1
\]
for the propagation of the perturbation $\delta\Xi$, by the standard
result for canonical scalars. This is
distinct from the \emph{adiabatic} sound speed
$c_{a}^{2} = \dot p_\Xi/\dot\rho_\Xi$, which depends on $w(z)$ and is
generally not equal to unity for a rolling quintessence field. The
quantity controlling clustering on subhorizon scales is $c_{s}^{2}$,
not $c_{a}^{2}$, by direct evaluation of \eqref{eq:rho}-\eqref{eq:p}
in the rest frame of the perturbation.
The model therefore predicts the standard quintessence behavior: dark-energy
perturbations are sound-speed-suppressed on subhorizon scales, and the
field does not cluster appreciably on cosmological scales of present
observational interest. On scales $k \gg a H$ the perturbations $\delta\Xi$
oscillate with frequency $c_s k/a = k/a$ and are damped relative to matter
perturbations; the contribution of $\delta\Xi$ to the late-time matter
power spectrum at $k \sim 0.1$~$h$/Mpc is therefore at the percent level
and below current Stage-IV sensitivity for distinguishing $c_s^{2} = 1$
quintessence from $\Lambda$CDM at the perturbation level
\cite{ShajibFrieman2025,Turyshev2026Review}. A full perturbation
calculation including ISW and CMB lensing predictions, and any departures
from $c_s^{2} = 1$ that arise from non-minimal couplings, is deferred to
the companion numerical paper.

\subsection{Gravitational-wave speed}

In \eqref{eq:action}, tensor perturbations propagate at the speed of
light, in accord with GW170817 \cite{GW170817}. No superluminal or
birefringent signal is predicted.

\section{Comparison to Literature}\label{sec:comparison}

Standard quintessence potentials, their existence of a minimum, their
late-time EoS behavior, and their comparison to the present model are
summarized in Table~\ref{tab:compare}.

\begin{table}[h]
\small
\begin{tabular}{l|c|c|c|l}
Potential & Reference & Minimum? & $w \to$ at late time & Notes\\
\hline
$V \sim \phi^{-\alpha}$ & \cite{RatraPeebles1988} & no & $w^{*} > -1$ const & tracker\\
$V \sim e^{-\alpha\phi}$ & \cite{Wetterich1988} & no & $w^{*}$ const & rolling\\
$V \sim 1 + \cos(\phi/f)$ & \cite{Frieman1995} & yes ($\phi = \pi f$) & approximate $-1$ & PNGB / thawing\\
$V_\text{CW}(\phi)$ & \cite{ColemanWeinberg1973} & yes, $\langle\phi\rangle(\lambda)$ & --- & loop-level $\phi^{4}\log\phi$, distinct family\\
$V \sim \phi^{p}(\log\phi)^{q}$ & \cite{BarrowParsons1995} & depends on $(p,q)$ & --- & inflation family\\
$\boxed{V = \Lambda^{4}\Xi\log\Xi}$ & present work & yes ($\Xi_0 = e^{-1}$) & $-1$ asymptotically & dual-regime freeze-thaw, $V = -\Lambda^{4} H_{\Gibbs}$
\end{tabular}
\caption{Comparison of quintessence potentials. The last row is the
model of the present paper. The closest structural relative is the
Barrow-Parsons inflation family \cite{BarrowParsons1995}
$V \sim V_0 \phi^{p}(\log\phi)^{q}$, which contains the present form as
the $(p,q) = (1,1)$ subcase but is applied to inflation, not dark energy,
and does not single out this subcase or derive its vacuum at $e^{-1}$.
Within the Caldwell-Linder \cite{CaldwellLinder2005} freezing/thawing
classification, the present trajectory's late-time asymptote is
freezing (Type-A, $w \to -1$); however on the broader trajectory the
field thaws outbound from near the vacuum and refreezes after a
Type-F turnaround at intermediate redshift $z_\star$ (\S\ref{sec:typeTFA}).
The single-history dual-regime traversal is not reproduced by any
single-regime quintessence in the table.}
\label{tab:compare}
\end{table}

\begin{remark}[On novelty]
Logarithmic functional forms appear in the dark-sector literature in
several distinct contexts. Thompson \cite{Thompson2019} studied
logarithmic and exponential dark-energy potentials in a beta-function
quintessence reconstruction framework, with logarithmic potentials
arising from the beta-function formalism rather than as fundamental
quintessence inputs. Kouwn and Oh \cite{KouwnOh2012} proposed a
dark-energy fluid with logarithmic equation of state. Oikonomou
\cite{Oikonomou2019} analyzed coupled dark-sector systems with
generalized logarithmic equations of state.

The most directly adjacent prior art is the logotropic-dark-fluid
program of Chavanis
\cite{Chavanis2015,Chavanis2021Kessence,Chavanis2022Logotropic}, which
proposes a unified dark fluid with a logarithmic equation of state
$P = A \ln(\rho/\rho_P)$ and, in its most recent
formulation~\cite{Chavanis2022Logotropic}, a relativistic complex
scalar field $\varphi$ with logarithmic potential $V(|\varphi|^{2})$
satisfying a generalized Klein-Gordon equation, with an explicit
classical limit $\hbar \to 0$ returning $\Lambda$CDM. Three structural
distinctions separate the present minimal theory from
\cite{Chavanis2022Logotropic} (and from the broader logotropic
program). \emph{First}, the field content here is a real, positive,
dimensionless scalar $\Xi$ rather than a complex scalar $\varphi$;
the action \eqref{eq:action} contains no $U(1)$ phase symmetry, no
charged-current sector, and no analogue of the Thomas-Fermi
approximation that the logotropic K-essence formalism
\cite{Chavanis2021Kessence} relies on. \emph{Second}, the present
work proposes $\Xi$ as a dark-energy field only and makes no claim of
dark-matter unification; the logotropic program is built explicitly
around DM/DE unification and is calibrated against the universal
surface density of dark-matter halos rather than against $w(z)$
evolution. \emph{Third}, and most directly observationally relevant,
the present model exhibits a dual-regime freeze-thaw transit
trajectory (\S\ref{sec:typeTFA}) on the rolling branch with outbound
initial condition, with a Type-F frozen turnaround at intermediate
redshift $z_\star$ and a non-monotone $w_{\rm DE}(z)$ signature
(F$_5$, \S\ref{sec:falsify}); the logotropic models with their
quasi-$\Lambda$CDM cosmological-scale behavior do not single out this
trajectory class. Other adjacent precedents using logarithmic
nonlinearity in nonlinear quantum mechanics
\cite{Zloshchastiev2010LogTime,Zloshchastiev2011SSB,AvdeenkovZloshchastiev2011,Zloshchastiev2018Stability,Zloshchastiev2020SuperfluidDM}
treat the nonlinearity as a quantum-wave-equation correction rather
than as a classical-scalar potential term, and are conceptually
distinct from the present minimal theory; we discuss this lineage
explicitly in \S\ref{sec:loglineage}.

The present work is distinct from each of these in three respects: (i)
it studies a specific dimensionless scalar with quintessence potential
$V = \Lambda^{4}\,\Xi\log\Xi$, (ii) it derives an exact analytic vacuum
at $\Xi_0 = e^{-1}$ that is independent of the coupling, and (iii) it
exhibits a dual-regime freeze-thaw transit trajectory --- Type-T
(thawing) $\to$ Type-F (frozen turnaround at $z_\star$) $\to$ Type-A
(asymptotic refreeze) --- within a single physical history, with the
non-monotone $w_{\rm DE}(z)$ signature predicted at intermediate
redshift (\S\ref{sec:typeTFA}, F$_5$). To the best of our
literature search
(covering arXiv categories astro-ph.CO, hep-ph, and gr-qc through April
2026, including DR2-era $V(\phi)$ reconstructions and analytic-family
surveys
\cite{AdamEtAl2025,Adil2026Analytic,Wang2026QReconstruction}), the
specific combination of features --- dimensionless field, the exact
$\Xi\log\Xi$ functional form, the coupling-independent $e^{-1}$ vacuum,
and the dual-regime freeze-thaw transit with explicit Type-F turnaround
at $z_\star \approx 2$ --- has not been previously studied as a
quintessence model. The Caldwell-Linder \cite{CaldwellLinder2005}
freezing/thawing classification covers monotone trajectories only;
the dual-regime traversal of the present model occupies a region of
quintessence parameter space outside that dichotomy. Closer to the
non-monotone class, work on emergent cosmologies with bouncing or
turning-point scalar dynamics has been explored
\cite{Bag2014NonMonotonic} but not within a logarithmic-potential
quintessence setting with the present trajectory taxonomy.
\end{remark}

\section{Scope and Limitations}\label{sec:scope}

The present paper asserts the following, and only the following:
\begin{itemize}
  \item The action \eqref{eq:action} is a minimal and self-consistent
        scalar-tensor theory. Its Euler-Lagrange equation \eqref{eq:EL},
        stress-energy tensor \eqref{eq:stress}, vacuum \eqref{eq:vacuum},
        fluctuation mass \eqref{eq:mass}, and FRW equations
        \eqref{eq:friedmann}-\eqref{eq:xieom} are formal consequences.
  \item The vacuum $\Xi_0 = e^{-1}$ coincides with the maximizer of the
        per-bin Gibbs integrand $-x\log x$ on $(0,1]$. As detailed in
        \S\ref{sec:entropy}, this is a structural relation between
        the functional form of $V$ and that of $H_{\Gibbs}$, not an
        identification of $\Xi$ with a probability density or of
        $-V/\Lambda^{4}$ with a thermodynamic entropy.
  \item The dual-regime freeze-thaw transit interpretation
        (\S\ref{sec:typeTFA}) is a property of the rolling branch with
        outbound initial condition $\dot\Xi_i > 0$ at $z_i \approx 20$.
        The Type-T/F/A classification organizes the trajectory's
        kinematics on this branch; the Type-F turnaround at
        $z_\star \approx 2$ is established empirically in the
        documented fit \eqref{eq:desibestfit} and depends on the
        choice of $(\Lambda^4, \Xi_i, \dot\Xi_i)$. Configurations with
        $\dot\Xi_i \leq 0$ realize different trajectory classes (pure
        Type-A asymptotic refreeze without an intermediate Type-F
        turnaround on the observable redshift interval); those
        configurations are not the subject of the present paper's
        empirical content.
  \item The DESI consistency check reported in \S\ref{sec:desifit} is a
        numerical integration of \eqref{eq:friedmann}-\eqref{eq:xieom}
        against the published DESI DR1 (2024) $(w_0, w_a)$ central
        values; the best-fit parameters \eqref{eq:desibestfit} are not
        the result of a full likelihood analysis including BAO, CMB,
        and SN data jointly, which is deferred to a companion numerical
        paper.
\end{itemize}

The present paper does \emph{not} claim the following:
\begin{itemize}
  \item Any derivation of $\Lambda$ from a microphysical substrate; the
        observation that $\Lambda$ is numerically near the dark-energy
        scale when $m_\Xi \sim H_0$ is a consistency check, not a
        solution to the cosmological-constant hierarchy problem.
  \item Any direct-matter coupling for $\Xi$, any observable fifth-force
        prediction, or any laboratory-scale consequence.
  \item Any mathematical uniqueness theorem for $V = \Xi\log\Xi$
        derived within the present setting. The Bialynicki-Birula--Mycielski
        uniqueness theorem \cite{BialynickiBirulaMycielski1976} is cited
        as structural motivation, not as a proof of uniqueness under
        the cosmological hypotheses.
  \item Any derivation of the cosmological action \eqref{eq:action} from
        the algebraic-combinatorial structure studied in the companion
        papers \cite{SandersGishSigma2026,SandersGishFourCore2026}. The
        companions establish a substrate non-associativity-decay rate
        and a closed-form algebraic attractor on a specific
        composition-table family at $N = 10$; the structural relations
        between that substrate and the logarithmic potential studied
        here are motivational, and the empirical content of the present
        paper is independent of those connections.
  \item Any analysis of the contribution of $\Xi$ during the
        inflationary epoch; the slow-roll parameters during inflation
        and any consequent effects on the primordial power spectrum are
        outside the scope of the present late-time dark-energy analysis.
\end{itemize}

\section{Summary}

A real, positive, dimensionless scalar field $\Xi$ with self-interaction
$V = \Lambda^{4}\,\Xi\log\Xi$, minimally coupled to gravity through
\eqref{eq:action}, defines a minimal dark-energy model with:
\begin{enumerate}
  \item Exact analytic vacuum $\Xi_0 = e^{-1}$
        (coupling-independent, universal).
  \item Stable massive fluctuations $m_\Xi^{2} = \Lambda^{4} e / M_\Pl^{2}$,
        with $\Lambda$ numerically near the observed dark-energy scale
        when $m_\Xi \sim H_0$.
  \item Freezing quintessence asymptotic behavior with $w \to -1$ in
        the cosmological future. On the dual-regime physical rolling
        branch (\S\ref{sec:typeTFA}), the trajectory thaws outbound,
        reaches a Type-F frozen turnaround at intermediate redshift
        $z_\star \approx 2$ where $\dot\Xi = 0$ and $w_\Xi(z_\star) = -1$
        instantaneously, and refreezes back toward the vacuum
        ($w_{\rm DE}(z)$ has a local minimum at $z_\star$).
  \item Structural relation $V = -\Lambda^{4}\,H_{\Gibbs}$ between the
        potential and the per-bin Gibbs integrand: the vacuum
        $\Xi_0 = e^{-1}$ coincides with the maximizer of $H_{\Gibbs}$
        on $(0,1]$. The connection is structural, not operational; see
        \S\ref{sec:entropy}.
  \item No observable fifth force is expected in the minimal theory.
  \item Consistency with DESI DR1 CPL central values at
        $\chi^{2}_{\text{Gauss}} = 1.52$ (1.2$\sigma$, compared to
        $\chi^{2}_{\text{Gauss}} = 15.26$ for $\Lambda$CDM under the
        same uncorrelated-Gaussian approximation on $(w_0,w_a)$) for
        $\Lambda \approx 1.7$ meV and reasonable initial conditions.
        These $\chi^{2}_{\text{Gauss}}$ values are not derived from the
        joint BAO + CMB + SN likelihood; full DR2 + SN + CMB likelihood
        analysis is deferred to a companion paper.
\end{enumerate}
The theory is falsifiable: a robust preference for phantom behavior on
the rolling branch, a non-monotone $w(z)$ trajectory on the inbound leg
($0 \leq z \leq 1.5$), failure of the asymptotic approach to $-1$, or
a confirmed monotonic $w_{\rm DE}(z)$ on $0 \leq z \leq 5$ with no local
minimum near $z_\star \approx 2$, would each rule out the corresponding
prediction of the minimal model at the precision of Stage-IV surveys.
The dual-regime signature (F$_5$, \S\ref{sec:falsify}) requires
non-parametric $w(z)$ reconstructions beyond the CPL parametrization,
since CPL is monotone by construction.

\section*{Acknowledgments}

The authors thank colleagues for discussions of the model's structural
properties. Verification scripts (\texttt{proof\_xi\_canonical.py}, 22
tests, all passing) and the DESI integration scripts
(\texttt{desi\_xi\_fit.py}, \texttt{desi\_xi\_optimize\_v2.py}) are available
as electronic supplementary material. All errors are the authors'.

\begin{thebibliography}{99}

\bibitem{BialynickiBirulaMycielski1976}
I.~Bialynicki-Birula and J.~Mycielski,
\emph{Nonlinear wave mechanics},
Ann. Phys. (N.Y.) \textbf{100}, 62--93 (1976).
DOI: 10.1016/0003-4916(76)90057-9.

\bibitem{CaldwellLinder2005}
R.~R.~Caldwell and E.~V.~Linder,
\emph{The limits of quintessence},
Phys. Rev. Lett. \textbf{95}, 141301 (2005); arXiv:astro-ph/0505494.

\bibitem{SandersGishSigma2026}
B.~R.~Sanders and M.~Gish,
\emph{Non-Associativity Decay in Binary Composition Tables over
$\mathbb{Z}/N\mathbb{Z}$},
Preprint, 2026; companion paper, submitted to
\emph{Journal of Combinatorial Theory, Series A};
verification scripts at DOI 10.5281/zenodo.18852047.

\bibitem{SandersGishFourCore2026}
B.~R.~Sanders and M.~Gish,
\emph{Joint Closure, Per-Coordinate Fuse Data, and a Closed-Form
Algebraic Attractor of Two Commutative Binary Operations on
$\mathbb{Z}/10\mathbb{Z}$},
Preprint, 2026; companion paper, submitted to
\emph{Algebraic Combinatorics};
verification scripts at DOI 10.5281/zenodo.18852047.

\bibitem{RatraPeebles1988}
B.~Ratra and P.~J.~E.~Peebles,
\emph{Cosmological consequences of a rolling homogeneous scalar field},
Phys. Rev. D \textbf{37}, 3406 (1988).

\bibitem{Wetterich1988}
C.~Wetterich,
\emph{Cosmology and the fate of dilatation symmetry},
Nucl. Phys. B \textbf{302}, 668 (1988).

\bibitem{Frieman1995}
J.~A.~Frieman, C.~T.~Hill, A.~Stebbins, and I.~Waga,
\emph{Cosmology with ultralight pseudo-Nambu-Goldstone bosons},
Phys. Rev. Lett. \textbf{75}, 2077 (1995).

\bibitem{CaldwellDaveSteinhardt1998}
R.~R.~Caldwell, R.~Dave, and P.~J.~Steinhardt,
\emph{Cosmological imprint of an energy component with general equation of state},
Phys. Rev. Lett. \textbf{80}, 1582 (1998); arXiv:astro-ph/9708069.

\bibitem{ZlatevWangSteinhardt1999}
I.~Zlatev, L.~Wang, and P.~J.~Steinhardt,
\emph{Quintessence, cosmic coincidence, and the cosmological constant},
Phys. Rev. Lett. \textbf{82}, 896 (1999); arXiv:astro-ph/9807002.

\bibitem{SteinhardtWangZlatev1999}
P.~J.~Steinhardt, L.~Wang, and I.~Zlatev,
\emph{Cosmological tracking solutions},
Phys. Rev. D \textbf{59}, 123504 (1999); arXiv:astro-ph/9812313.

\bibitem{SahniStarobinsky2000}
V.~Sahni and A.~A.~Starobinsky,
\emph{The case for a positive cosmological $\Lambda$-term},
Int. J. Mod. Phys. D \textbf{9}, 373 (2000); arXiv:astro-ph/9904398.

\bibitem{PeeblesRatra2003}
P.~J.~E.~Peebles and B.~Ratra,
\emph{The cosmological constant and dark energy},
Rev. Mod. Phys. \textbf{75}, 559 (2003); arXiv:astro-ph/0207347.

\bibitem{Padmanabhan2003}
T.~Padmanabhan,
\emph{Cosmological constant: the weight of the vacuum},
Phys. Rep. \textbf{380}, 235 (2003); arXiv:hep-th/0212290.

\bibitem{CopelandSamiTsujikawa2006}
E.~J.~Copeland, M.~Sami, and S.~Tsujikawa,
\emph{Dynamics of dark energy},
Int. J. Mod. Phys. D \textbf{15}, 1753 (2006); arXiv:hep-th/0603057.

\bibitem{FriemanTurnerHuterer2008}
J.~A.~Frieman, M.~S.~Turner, and D.~Huterer,
\emph{Dark energy and the accelerating universe},
Annu. Rev. Astron. Astrophys. \textbf{46}, 385 (2008); arXiv:0803.0982.

\bibitem{Tsujikawa2013}
S.~Tsujikawa,
\emph{Quintessence: a review},
Class. Quantum Gravity \textbf{30}, 214003 (2013); arXiv:1304.1961.

\bibitem{JoyceJainKhouryTrodden2015}
A.~Joyce, B.~Jain, J.~Khoury, and M.~Trodden,
\emph{Beyond the cosmological standard model},
Phys. Rep. \textbf{568}, 1 (2015); arXiv:1407.0059.

\bibitem{CPL2001}
M.~Chevallier and D.~Polarski,
\emph{Accelerating universes with scaling dark matter},
Int. J. Mod. Phys. D \textbf{10}, 213 (2001).

\bibitem{Linder2003}
E.~V.~Linder,
\emph{Exploring the expansion history of the universe},
Phys. Rev. Lett. \textbf{90}, 091301 (2003); arXiv:astro-ph/0208512.

\bibitem{JassalBaglaPadmanabhan2005}
H.~K.~Jassal, J.~S.~Bagla, and T.~Padmanabhan,
\emph{WMAP constraints on low redshift evolution of dark energy},
Mon. Not. Roy. Astron. Soc. \textbf{356}, L11 (2005); arXiv:astro-ph/0404378.

\bibitem{BarbozaAlcaniz2008}
E.~M.~Barboza Jr. and J.~S.~Alcaniz,
\emph{A parametric model for dark energy},
Phys. Lett. B \textbf{666}, 415 (2008); arXiv:0805.1713.

\bibitem{Linder2006Paths}
E.~V.~Linder,
\emph{Paths of quintessence},
Phys. Rev. D \textbf{73}, 063010 (2006); arXiv:astro-ph/0601052.

\bibitem{ScherrerSen2008Thawing}
R.~J.~Scherrer and A.~A.~Sen,
\emph{Thawing quintessence with a nearly flat potential},
Phys. Rev. D \textbf{77}, 083515 (2008); arXiv:0712.3450.

\bibitem{CrittendenMajerottoPiazza2007}
R.~Crittenden, E.~Majerotto, and F.~Piazza,
\emph{Measuring deviations from a cosmological constant: a field-space parametrization},
Phys. Rev. Lett. \textbf{98}, 251301 (2007); arXiv:astro-ph/0702003.

\bibitem{CrittendenPogosianZhao2009}
R.~G.~Crittenden, L.~Pogosian, and G.-B.~Zhao,
\emph{Investigating dark energy experiments with principal components},
JCAP \textbf{12}, 025 (2009); arXiv:astro-ph/0510293.

\bibitem{Holsclaw2010GP}
T.~Holsclaw, U.~Alam, B.~Sansó, H.~Lee, K.~Heitmann, S.~Habib, and D.~Higdon,
\emph{Nonparametric dark energy reconstruction from supernova data},
Phys. Rev. Lett. \textbf{105}, 241302 (2010); arXiv:1011.3079.

\bibitem{SahniShafielooStarobinsky2008}
V.~Sahni, A.~Shafieloo, and A.~A.~Starobinsky,
\emph{Two new diagnostics of dark energy},
Phys. Rev. D \textbf{78}, 103502 (2008); arXiv:0807.3548.

\bibitem{PogosianShafieloo2009}
L.~Pogosian and A.~Shafieloo,
\emph{Tomographic reconstruction of dark energy},
Phys. Rev. D \textbf{84}, 023523 (2011); arXiv:0904.1029.

\bibitem{Albrecht2006DETF}
A.~Albrecht \emph{et al.},
\emph{Report of the Dark Energy Task Force},
arXiv:astro-ph/0609591 (2006).

\bibitem{Aubourg2015}
\'{E}.~Aubourg \emph{et al.},
\emph{Cosmological implications of baryon acoustic oscillation measurements},
Phys. Rev. D \textbf{92}, 123516 (2015); arXiv:1411.1074.

\bibitem{EuclidLaureijs2011}
R.~Laureijs \emph{et al.} (Euclid Collaboration),
\emph{Euclid Definition Study Report},
arXiv:1110.3193 (2011).

\bibitem{LSSTIvezic2019}
\v{Z}.~Ivezi\'{c} \emph{et al.} (LSST Collaboration),
\emph{LSST: from science drivers to reference design and anticipated data products},
Astrophys. J. \textbf{873}, 111 (2019); arXiv:0805.2366.

\bibitem{BarrowParsons1995}
J.~D.~Barrow and P.~Parsons,
\emph{Inflationary models with logarithmic potentials},
Phys. Rev. D \textbf{52}, 5576 (1995); arXiv:astro-ph/9506049.

\bibitem{ColemanWeinberg1973}
S.~Coleman and E.~Weinberg,
\emph{Radiative corrections as the origin of spontaneous symmetry breaking},
Phys. Rev. D \textbf{7}, 1888 (1973).

\bibitem{Rosen1969}
G.~Rosen,
\emph{Dilatation covariance and exact solutions in local relativistic field theories},
Phys. Rev. \textbf{183}, 1186 (1969).

\bibitem{CazenaveHaraux1980}
T.~Cazenave and A.~Haraux,
\emph{Equations d'\'evolution avec non lin\'earit\'e logarithmique},
Ann. Fac. Sci. Toulouse \textbf{2}, 21--51 (1980).

\bibitem{HoeghKrohn1971}
R.~H\o egh-Krohn,
\emph{A general class of quantum fields without cut-offs in two space-time dimensions},
Commun. Math. Phys. \textbf{21}, 244--255 (1971).

\bibitem{Hefter1985}
E.~F.~Hefter,
\emph{Application of the nonlinear Schr\"{o}dinger equation with a logarithmic
inhomogeneous term to nuclear physics},
Phys. Rev. A \textbf{32}, 1201 (1985).

\bibitem{Weinberg1989Precision}
S.~Weinberg,
\emph{Precision tests of quantum mechanics},
Phys. Rev. Lett. \textbf{62}, 485 (1989).

\bibitem{Bollinger1989}
J.~J.~Bollinger, D.~J.~Heinzen, W.~M.~Itano, S.~L.~Gilbert, and D.~J.~Wineland,
\emph{Test of the linearity of quantum mechanics by rf spectroscopy of the
$^{9}$Be$^{+}$ ground state},
Phys. Rev. Lett. \textbf{63}, 1031 (1989).

\bibitem{Zloshchastiev2010LogTime}
K.~G.~Zloshchastiev,
\emph{Logarithmic nonlinearity in theories of quantum gravity: origin of time
and observational consequences},
Gravit. Cosmol. \textbf{16}, 288--297 (2010); arXiv:0906.4282.

\bibitem{AvdeenkovZloshchastiev2011}
A.~V.~Avdeenkov and K.~G.~Zloshchastiev,
\emph{Quantum Bose liquids with logarithmic nonlinearity: self-sustainability
and emergence of spatial extent},
J. Phys. B \textbf{44}, 195303 (2011); arXiv:1108.0847.

\bibitem{Zloshchastiev2011SSB}
K.~G.~Zloshchastiev,
\emph{Spontaneous symmetry breaking and mass generation as built-in phenomena
in logarithmic nonlinear quantum theory},
Acta Phys. Polon. B \textbf{42}, 261 (2011); arXiv:0912.4139.

\bibitem{Zloshchastiev2018Stability}
K.~G.~Zloshchastiev,
\emph{Stability and metastability of trapless Bose-Einstein condensates and
quantum liquids},
Z. Naturforsch. A \textbf{72}, 677 (2017); arXiv:1709.06694.

\bibitem{Zloshchastiev2020SuperfluidDM}
K.~G.~Zloshchastiev,
\emph{An alternative to dark matter and dark energy: scale-dependent gravity
in superfluid vacuum theory},
Universe \textbf{6}, 180 (2020); arXiv:2009.03126.

\bibitem{Chavanis2015}
P.-H.~Chavanis,
\emph{The logotropic dark fluid as a unification of dark matter and dark energy},
Phys. Rev. D \textbf{92}, 103004 (2015); arXiv:1505.00034.

\bibitem{Chavanis2022Logotropic}
P.-H.~Chavanis,
\emph{New logotropic model based on a complex scalar field with a logarithmic potential},
Phys. Rev. D \textbf{106}, 063525 (2022); arXiv:2201.05908.

\bibitem{Chavanis2021Kessence}
P.-H.~Chavanis,
\emph{$K$-essence Lagrangians of polytropic and logotropic unified
dark matter and dark energy models},
Astronomy \textbf{1}, 126 (2022); arXiv:2109.05963.

\bibitem{Caldwell2002Phantom}
R.~R.~Caldwell,
\emph{A phantom menace? Cosmological consequences of a dark energy
component with super-negative equation of state},
Phys. Lett. B \textbf{545}, 23 (2002); arXiv:astro-ph/9908168.

\bibitem{Bag2014NonMonotonic}
S.~Bag, V.~Sahni, Y.~Shtanov, and S.~Unnikrishnan,
\emph{Emergent cosmology revisited},
JCAP \textbf{07}, 034 (2014); arXiv:1403.4243.

\bibitem{Planck2020}
Planck Collaboration,
\emph{Planck 2018 results. VI. Cosmological parameters},
Astron. Astrophys. \textbf{641}, A6 (2020).

\bibitem{DESI2024BAO}
DESI Collaboration,
\emph{DESI 2024 III: Baryon acoustic oscillations from galaxies and quasars},
arXiv:2404.03000 (2024).

\bibitem{DESI2024VI}
DESI Collaboration,
\emph{DESI 2024 VI: Cosmological constraints from the measurements of baryon acoustic oscillations},
arXiv:2404.03002 (2024).

\bibitem{DESI2025DR2}
DESI Collaboration,
\emph{DESI DR2 results. II. Measurements of baryon acoustic oscillations and cosmological constraints},
arXiv:2503.14738; Phys. Rev. D \textbf{112}, 083515 (2025).

\bibitem{Wang2026QReconstruction}
S.~Wang, T.-N.~Li, T.~Liu, and G.-H.~Du,
\emph{Model-independent reconstruction of quintessence potential and kinetic energy from DESI DR2 and Pantheon+ supernovae},
arXiv:2603.21125 (2026).

\bibitem{Adil2026Analytic}
S.~A.~Adil, M.~A.~Zapata, \"{O}.~Akarsu, and J.~A.~Vazquez,
\emph{Background-level reconstruction of scalar-field potentials from dark-energy histories and comparison with analytic potential families},
arXiv:2603.14693 (2026).

\bibitem{AdamEtAl2025}
H.~Adam, M.~P.~Hertzberg, D.~Jim\'enez-Aguilar, and I.~Khan,
\emph{Comparing Minimal and Non-Minimal Quintessence Models to 2025 DESI Data},
arXiv:2509.13302 (2025).

\bibitem{Turyshev2026Review}
S.~G.~Turyshev,
\emph{Dark Energy After DESI DR2: Observational Status, Reconstructions, and Physical Models},
arXiv:2602.05368 (2026).

\bibitem{ShajibFrieman2025}
A.~J.~Shajib and J.~A.~Frieman,
\emph{Scalar field dark energy models: current and forecast constraints},
arXiv:2502.06929; Phys. Rev. D \textbf{112}, 063508 (2025).

\bibitem{Thompson2019}
R.~I.~Thompson,
\emph{Beta function quintessence cosmological parameters and fundamental
constants -- II. Exponential and logarithmic dark energy potentials},
Mon. Not. Roy. Astron. Soc. \textbf{482}, 5448 (2019); arXiv:1811.03164.

\bibitem{KouwnOh2012}
S.~Kouwn and P.~Oh,
\emph{Dark energy with logarithmic cosmological fluid},
J. Korean Phys. Soc. \textbf{65}, 814 (2014); arXiv:1201.4544.

\bibitem{Oikonomou2019}
V.~K.~Oikonomou,
\emph{Generalized logarithmic equation of state in classical and loop
quantum cosmology dark energy--dark matter coupled systems},
Annals of Physics \textbf{409}, 167912 (2019); arXiv:1907.02600.

\bibitem{GW170817}
B.~P.~Abbott \emph{et al.} (LIGO Scientific Collaboration and Virgo Collaboration),
\emph{Gravitational waves and gamma-rays from a binary neutron star merger:
GW170817 and GRB 170817A},
Astrophys. J. Lett. \textbf{848}, L13 (2017).

\bibitem{Verlinde2011}
E.~Verlinde,
\emph{On the origin of gravity and the laws of Newton},
JHEP \textbf{04}, 029 (2011).

\bibitem{CHLZ2012}
A.~Sheykhi \emph{et al.},
\emph{Entropy-corrected holographic dark energy and interacting quintessence reconstruction},
Astrophys. Space Sci. \textbf{342}, 93 (2012).

\bibitem{Ensslin2013}
T.~A.~En{\ss}lin,
\emph{Information theory for fields},
Annalen der Physik \textbf{525}, 895 (2013); arXiv:1301.2556.

\bibitem{Caticha2012}
A.~Caticha,
\emph{Entropic inference and the foundations of physics},
EBEB 2012 monograph (2012); arXiv:1412.5629.

\bibitem{EotWash}
D.~J.~Kapner \emph{et al.},
\emph{Tests of the gravitational inverse-square law below the dark-energy length scale},
Phys. Rev. Lett. \textbf{98}, 021101 (2007).

\bibitem{MICROSCOPE}
P.~Touboul \emph{et al.},
\emph{MICROSCOPE mission: final results of the test of the equivalence principle},
Phys. Rev. Lett. \textbf{129}, 121102 (2022).

\end{thebibliography}

\end{document}